% ID: SEL-D0054-OB-E01
% ESTADO: validada. S055.
\textbf{a)} Aunque $4{,}1+11{,}7=15{,}8\le20$, se tiene $3\cdot4{,}1+5\cdot11{,}7=70{,}8>70$. Por tanto, el punto queda fuera.

\textbf{b)} En los ejes, las cotas más restrictivas son $x=20$ e $y=14$. Las dos rectas oblicuas se cortan en $(15,5)$, como resulta al sustituir $x=20-y$ en $3x+5y=70$. Junto con el origen, son los cuatro vértices.
\begin{center}
\setlength{\unitlength}{0.9mm}
\begin{picture}(86,56)(-8,-8)
\qbezier(0,0)(34,0)(68,0)
\qbezier(0,0)(0,22)(0,44)
\put(69,0){$x$}
\put(2,43){$y$}
\qbezier(60,0)(52.5,7.5)(45,15)
\qbezier(45,15)(22.5,28.5)(0,42)
\put(53,-6){\scriptsize $(20,0)$}
\put(46,16){\scriptsize $(15,5)$}
\put(2,37){\scriptsize $(0,14)$}
\put(-6,-6){\scriptsize $(0,0)$}
\end{picture}
\end{center}
\textbf{c)} El mínimo es $\boxed{0}$, sólo en el origen, pues los coeficientes de $F$ son positivos y $x,y\ge0$. Como $5F=3x+5y\le70$, el máximo es $\boxed{14}$ en todo el segmento entre $(0,14)$ y $(15,5)$, que contiene exactamente los puntos factibles donde hay igualdad.
